% ---- Appendix F: add after the paragraph on the naive baseline -------------------
The naïve baseline is the model's mean accuracy over the training languages. To isolate what the
language features and interactions add, we also compare against two simpler predictors
(Table~\ref{tab:lolo-baselines}): a size-only model, and a binomial regression on model identity and
log Common Crawl pages (the resource proxy of Equation~1), refitted with each language held out.
Model identity and resource level alone recover most of the gain over the naïve baseline (language
mean 8.5$\to$3.4 pp, model in language 9.8$\to$6.5 pp). The full model improves on this by a further
0.3 and 0.5 pp (9\% and 8\%), but with 15 languages the difference is not distinguishable from zero
(95\% bootstrap intervals $[-1.4, 0.8]$ and $[-1.3, 0.3]$ pp). The predictive gain over the naïve
baseline is therefore attributable mainly to resource level; the remaining features and interactions
add explanatory structure (Section~4.1) without a measurable gain in out-of-language forecasting.
Evaluating a few templates in the target language reduces the error further (Table~\ref{tab:fewshot}).

% ---- Main text, Finding 5 (Section 4.2): replace the numbers ---------------------
A language's average difficulty is predictable to 3.1 pp (Pearson $r = 0.93$, Fig.~5) and, for a
specific model in that language, to 6.1 pp ($r = 0.96$, Fig.~6), a 64\% and 38\% reduction over the
naïve baseline. Much of this is carried by language resource level: a predictor using only model
identity and resource level reaches 3.4 pp and 6.6 pp (Appendix~F), and the full model's further
improvement is not distinguishable from zero. Utilising just 10 templates further reduces the error to
2.0 pp and 4.4 pp (Table~\ref{tab:fewshot}).
